\chapter{30}

30 er et terningespil med 6 terninger, hvor det gælder om at slå over $\geq 30$ eller $\leq 10$.

Til spillet skal der bruges: 1 raflebæger og 6 terninger. Man starter med at slå med alle 6 terninger, og tage minimum én terning fra\footnote{Man SKAL tage én terning hver gang man laver et slag.}. Man slår da igen med de 5 (eller mindre) terninger og igen tager minimum 1 terning fra. Dette fortsættes indtil man ikke kan tage flere terninger fra og man summer alle terningernes øjne sammen. 
Hvis man har slået $<30$ så skal man drikke tåre op til $30$. F.eks. hvis der er en person der har slået $4 + 5 + 6 + 4 + 2 + 6 = 27$ så skal de drikke 3 tåre.
Hvis man har slået præcis $30$ så er der fælleskål og alle drikker.
Hvis man har slået $>30$ så skal man slå efter det tal man har over $30$. Så hvis man har slået $32$ skal man slå efter $2$'er, eller hvis man har slået $34$ skal man slå efter $4$'er.
Igen slår man med alle $6$ terninger og, hvis man slår det tal man er efter eksempelvis $2$'er så tager man dem fra og slår igen med de resterende terninger indtil man ikke slår flere $2$'er\footnote{Eller hvilket tal man nu er efter.}. Man summer igen terningernes øjne og det er de tåre man kan give ud, de skal dog gives ud multiplum af det man har slået. Så eksemplevis, hvis en person har slået tre $4$'er så kan den person give ud 4 tåre af tre omgange.

Når man skal slå efter tallet over $30$, og man slår med de $6$ terninger, så hvis man IKKE slår nogen af det tal man er efter, så slå man mod sig selv. Her er reglerne det samme som før, blot at du selv skal drikke alle de tåre du slår. Så i det ovenstående tilfælde ville man skulle drikke $12$ tåre selv.

Dette "frem og tilbage" mellem dig selv og dine modstandere laves indtil første gang der bliver slået det tal du er efter.

Derefter er det den næste persons tur og man spiller indtil alle er gået kolde eller man ikke gider mere.
